\section{Pole gluing, Robin response and coherent delays}
\label{sec:poles}
There are two explicit ways to produce the pole kernel
$P_L(x,y)=2\cosh((x-y)/2)$. The first changes a boundary response
and exposes an unstable auxiliary mode. The second interferes
with remote delayed sources and adds a positive contact.

\subsection{The attractive Robin identity}
Put $a=1/2$, $h=L/2$, and let $R_{\rm f}$ be the interval integral
operator with kernel $e^{-|x-y|/2}$. The lowest gamma contribution
is $K_0=4I-R_{\rm f}$. Define
\[
S_{\rm R}u=-u'',\qquad A_{\rm R}=S_{\rm R}+a^2,
\]
\[
\Dom S_{\rm R}=\{u\in H^2(-h,h):
 u'(h)=au(h),\ u'(-h)=-au(-h)\}.
\]
Its closed form on $H^1(-h,h)$ is
\[
\mathfrak a_{\rm R}[u]=
 \int_{-h}^h(|u'|^2+a^2|u|^2)\,dx
 -a(|u(h)|^2+|u(-h)|^2).
\]
Endpoint trace bounds prove semiboundedness and closedness. Direct
differentiation, the derivative jump and the Robin conditions give
\begin{equation}\label{poles:robin-green}
A_{\rm R}^{-1}(x,y)=-e^{|x-y|/2}
                  =R_{\rm f}(x,y)-P_L(x,y),
\qquad
C_L:=K_0+P_L=4I-A_{\rm R}^{-1}.
\end{equation}
There is no zero mode of $A_{\rm R}$.

The free whole-line auxiliary field, eliminated outside the interval,
has energy
$4\norm{f-u}^2+16\norm{u'}^2+8(|u(h)|^2+|u(-h)|^2)$.
Producing the pole law reverses the sign of the endpoint term.
Equivalently the resulting energy is
\[
\mathcal E_{\rm R}(f,u)=4\norm f^2-8\re\ip f u
                                      +16\mathfrak a_{\rm R}[u].
\]
Its stationary point is $u_f=A_{\rm R}^{-1}f/4$, and
\begin{equation}\label{poles:stationary}
\mathcal E_{\rm R}(f,u)=\ip f{C_Lf}
                    +16\mathfrak a_{\rm R}[u-u_f].
\end{equation}
This is a stationary elimination, not a nonnegative minimization.

Indeed $A_{\rm R}$ has exactly one negative eigenvalue $-\delta_e$,
where
\[
\kappa_e\tanh(\kappa_e h)=a,\qquad
\delta_e=\kappa_e^2-a^2>0,\qquad
\phi_e(x)\propto\cosh(\kappa_e x).
\]
The even equation has a unique root $\kappa_e>a$. An odd negative
eigenvalue of $A_{\rm R}$ would require
$\kappa\coth(\kappa h)=a$ with $\kappa>a$, which is impossible.
Thus $\inf_u\mathcal E_{\rm R}(f,u)=-\infty$ for every $f$.

\begin{proposition}\label{poles:threshold}
The bounded reduced operator $C_L$ is strictly positive for
$L<4$, has kernel spanned by $x$ at $L=4$, and has exactly one
negative eigenvalue, in the odd sector, for $L>4$.
\end{proposition}
\begin{proof}
On an eigenvector of $S_{\rm R}$ with eigenvalue $\nu$, the
coefficient of $C_L$ is $4\nu/(a^2+\nu)$.
The even negative eigenvalue has $\nu<-a^2$ and gives a positive
coefficient. An odd negative eigenvalue of $S_{\rm R}$ exists
exactly when $ah>1$, from
$\kappa_o\coth(\kappa_oh)=a$, $0<\kappa_o<a$.
It gives the single negative coefficient. At $ah=1$, the odd
zero mode is $x$. The remaining eigenvalues are positive.
As direct checks,
$C_L1=4e^{L/4}\cosh(x/2)$ and
$C_Lx=8e^{L/4}(1-L/4)\sinh(x/2)$.
\end{proof}
For $L>4$ write
$r_L=(a^2-\kappa_o^2)^{-1}-4>0$ and let $\Pi_o$ project onto the
normalized odd mode. Set $r_L\Pi_o=0$ for $L\le4$.
Removing the negative even auxiliary direction makes
\eqref{poles:stationary} a minimization, but its reduced operator
is $4I-A_{\rm R}^{-1}(I-\Pi_e)$; recovering $C_L$ adds
$\delta_e^{-1}\Pi_e$. This does not remove the odd reduced defect.

There is also an explicit positive source for $C_L+r_L\Pi_o$.
Let $S_+$ be the nonnegative part of $S_{\rm R}$ and $P_{\ge0}$
its spectral projection. Use
\[
\mathcal B_Lf=
\begin{pmatrix}
2S_+(S_++1/4)^{-1}P_{\ge0}f\\
S_+^{1/2}(S_++1/4)^{-1}P_{\ge0}f\\
\sqrt{4+\delta_e^{-1}}\,\ip{\phi_e}f
\end{pmatrix}.
\]
Functional calculus gives
$\mathcal B_L^*\mathcal B_L=C_L+r_L\Pi_o$.
All components are bounded; the adjoint is the sum of the
transposed components with conjugated scalar coefficients.
Adjoining the higher gamma modes yields a closed source of norm
$K+P_L+r_L\Pi_o$. The full arithmetic form still has the
contact, translations and $-r_L\Pi_o$ to account for.

\subsection{A supersymmetric factorization needs a shift}
With $w(x)=\kappa_e\tanh(\kappa_ex)$, take
$D=\partial_x-w$ on $H^1$ and
$D^*=-\partial_x-w$ on $H_0^1$. Then
\[
D^*D=S_{\rm R}+\kappa_e^2=A_{\rm R}+\delta_e,\qquad
DD^*=-\partial_x^2+\kappa_e^2
          -2\kappa_e^2\operatorname{sech}^2(\kappa_ex)
\]
with Dirichlet partner conditions. This is consistent with the
boundary-domain requirements in supersymmetric quantum mechanics
\cite{BoundarySUSY}. Using $D^*D$ as the kinetic operator in the
resolvent changes the reduced operator to
$C_{\rm shift}=4I-(A_{\rm R}+\kappa_e^2)^{-1}$.
Its difference from the desired one is
\[
C_L-C_{\rm shift}
 =-\kappa_e^2 A_{\rm R}^{-1}(A_{\rm R}+\kappa_e^2)^{-1}.
\]
It is infinite rank. Positive factorization of a shifted
Hamiltonian does not restore the unshifted pole identity.

\subsection{Remote interference with a gamma channel}
Fix a mass $a$ and put $\alpha_a=\sqrt{2/a}$. Its gamma column is
$\alpha_a((I-a^2R_a)F,a\partial_xR_aF)$.
Tensoring with the unit quartic cap $\chi$ in
Section~\ref{sec:caps} leaves all the following Gram identities
unchanged. Add to this same output channel
\[
\mathcal B_{D,t}F=-t
\begin{pmatrix}(U_D+U_{-D})F\\(U_D-U_{-D})F\end{pmatrix},
\qquad t\in\R.
\]
Its square is $4t^2I$. Writing $V=U_D+U_{-D}$ and
$J=U_D-U_{-D}$, its cross term with the gamma source is
\[
-2\alpha_atV+2\alpha_at a^2R_aV
                   +2\alpha_a at\,\partial_xR_aJ.
\]
The multiplier is
$-4\alpha_at(\tau^2\cos(D\tau)-a\tau\sin(D\tau))/(a^2+\tau^2)$.
For $D\ge L$ the atoms disappear after interval compression and
the remaining kernel is
$4\alpha_a at e^{-aD}\cosh(a(x-y))$.
Taking $a=1/2$ and $t=e^{D/2}/2$ gives a closed source $A_0$ with
\begin{equation}\label{poles:delay-source}
T_0:=A_0^*A_0=T_L+P_L+e^D I,\qquad
T_0\ge mI,\quad m=e^D+L-2\sinh(L/2)>0.
\end{equation}
The lower bound drops the nonnegative gamma and even pole
components and uses $\norm s^2=\sinh(L/2)-L/2$.
Here and below $D\ge L$ is fixed.

The contact $e^D$ is optimal in the two-remote-shift ansatz with
arbitrary constant internal cap vectors. If $r_\pm,s_\pm$ are
their scalar overlaps with the two gamma residuals, the pole
coefficient $\kappa\cosh(a(x-y))$ fixes
$X=r_++\bar r_-+s_+-\bar s_-=-\kappa e^{aD}/(\alpha_a a)$.
Cauchy--Schwarz bounds the contact below by
\[
|X|^2/4=\kappa^2e^{2aD}/(8a).
\]
For the pole normalization this is $e^D$, attained above.
The statement is restricted to this ansatz.

\subsection{Finite delay rigidity}
A one-sided prime injection
$-(c/\alpha_a)(U_dF,0)$ has square plus cross term
\[
\frac{ac^2}{2}I-c(U_d+U_{-d})
                   +ca^2R_a(U_d+U_{-d}).
\]
The final term is a continuous kernel
$ca(e^{-a|x-y-d|}+e^{-a|x-y+d|})/2$.
Thus an isolated shifted source produces the desired atoms along
with a continuous distortion and a contact. Several shifts also
produce autocorrelations between their outputs.

More generally let the delay alphabet $\mathcal D$ be finite and
let the cap vectors in channel $k$ and residual component
$\varepsilon=0,1$ obey
\[
\sum_{k,\varepsilon}
       \left(\sum_{d\in\mathcal D}\norm{z_{k,\varepsilon,d}}\right)^2
<\infty .
\]
This defines a bounded finite-delay source $\mathcal B$.
Its zero-delay autocorrelation coefficient is the nonnegative
$\gamma_{\mathcal B}=\sum_{k,\varepsilon,d}
\norm{z_{k,\varepsilon,d}}^2$.

\begin{proposition}\label{poles:finite-delay}
Within this class, if the continuous gamma kernel is preserved
(or changed only by the specified pole kernel), every added
interior atom comes from $\mathcal B^*\mathcal B$.
The contact coefficient is $\gamma_{\mathcal B}$.
The cross term with the gamma source cannot independently supply
nonzero interior prime atoms while its continuous part is
cancelled.
\end{proposition}
\begin{proof}
Let $u_{k,d},w_{k,d}$ be the overlaps of the cap vectors with the
two gamma residuals, and put
$\sigma_{k,d}=u_{k,d}+\bar u_{k,-d}$,
$\theta_{k,d}=\bar w_{k,-d}-w_{k,d}$.
The cross term in channel $k$ is
\[
\alpha_{a_k}\sum_d\sigma_{k,d}U_d
-\alpha_{a_k}a_k^2R_k\sum_d\sigma_{k,d}U_d
+\alpha_{a_k}a_k\partial_xR_k\sum_d\theta_{k,d}U_d.
\]
Away from the finitely many centers, its continuous kernel is
\[
-\sum_{k,d}\frac{\alpha_{a_k}a_k}{2}
 [\sigma_{k,d}+\theta_{k,d}\operatorname{sgn}(z-d)]
                              e^{-a_k|z-d|}.
\]
On each interval between centers this is an absolutely convergent
Laurent expansion in $\zeta=e^{z/2}$ with powers $\pm(4k+1)$.
Uniqueness of Laurent coefficients identifies it with zero, or
with the two prescribed pole powers. Comparing adjacent intervals
forces $\sigma_{k,d}=\theta_{k,d}=0$ at every interior center.
Remote centers are smooth across the interval and supply no
interior atoms. The source autocorrelation is therefore the only
remaining atomic contribution, with the stated zero coefficient.
\end{proof}
Filters, resolvent feedback, and infinite alphabets with accumulating
delays are outside this result. These distinctions matter in the
constructions below.
